\chapter{\label{ch:theory}Foundations}

\minitoc

In this thesis experimental studies of quantum transport in carbon nanostructures have been carried out.

More generally, quantum transport phenomena occur on a mesoscopic scale, that is, when dealing with condensed phase systems, in particular (but not exclusively) those of micro-nano dimensions, 
for which behaviours not present in macroscopic bodies are observed.

The different behaviour of mesoscopic phases compared to macroscopic ones lies in the wavelength nature of electrons. This can be well understood by considering the different length scales which take
place in transport phenomena. These includes the size of the system $L$ (which can be different for the three directions in space, $L_x, L_y, L_z$), the Fermi wavelength $\lambda_F = h / \sqrt{2m_eE_F} = 2\pi/k_F$, and the coherence length $L_\phi$, which 
is the spatial distance that an electron travels before the phase of its wave function changes of $2\pi$. 

Additionally, two characteristic lengths associated with the mean distance traveled after elastic ($l_\textrm{el}$) and inelastic ($l_\textrm{in}$) collisions are also introduced; for inelastic collision, it is convenient to consider the mean distance before losing an amount of energy $k_BT$.      

For a metal with carrier concentrations $n \sim 10^{23}$ cm$^{-3}$, $l = 1 - 10$ nm, $\lambda_F = 0.5$ nm, and $L_\phi = 0.5$ \si{\micro\metre}, where $l$ is the mean free path and includes both elastic and inelastic events.
In comparison, for carrier concentrations $n \sim 10^{11}$ cm$^{-2}$, graphene displays $l = 50$ nm$-3$\footnote{For suspended graphene, $l (\textrm{4 K}) \approx 100$ nm, while Mayorov et al. reported values up to 3 \si{\micro\metre} at 290 K for graphene sandwiched between two hBN crystals \cite{mayorov2011micrometer}}  \si{\micro\metre}, $\lambda_F = 2\sqrt{\pi/n}$ nm, and $L_\phi = 0.5$ \si{\micro\metre} \cite{torres2013introduction}.




% A solid sample can be considered macroscopic if its properties are independent of the size of the sample. If we proceed, however, to measure the electric current in nanoscale metallic samples, 
% we reach a dimension below which the properties become dependent on the sample. Similar behaviour occurs in general when a solid body is divided into parts of mesoscopic size. 
% In the mesoscopic regime on two samples of the same size and chemical composition, therefore apparently identical, different properties can be measured, that is different values of the same physical quantity, such as
% values of the same physical quantity, such as electric current or thermal conduction.
% These variations from sample to sample are called mesoscopic fluctuations. This implies that the definition of these properties is only possible by means of an ensemble average. In a macroscopic solid this average is not necessary.
% In mesoscopic physics particular interest has been directed to the properties of electric transport, in particular to the conduction in direct current in linear regime in micro-nano-sized devices.


\section{Charge Transport In Zero-Dimensional Objects}

quantum dots transport physics

\section{Charge Transport in One-Dimensional Objects}

nanowires transport physics

\section{Thermoelectric Phenomena in Zero-Dimensional Structures}

What are the necessary changes that one has to take into account
when dealing also with other extensive variable such as a temperature gradient

\section{Density Functional Theory}

What are the relevant parameters/quantities that we are interested to calculate ab initio
Why we opted for a pseudopotential such as SIESTA rather than others (quantum espresso).
Explain why calculating the isolated gas phase all-electron electronic structure is legit for a 
qualitative see how the electron density is distributed over the system.

Remember also to emphasise the difference from stricly and non-stricly localised functions
used for computing the electronic properties: when the functions are stricly localised they are zero
outside a given radius from the atom they are associated with and so the radius is a key quality determined parameter for the basis set


\section{Non-Equilibrium Green Function Theory}

This section is needed to explain formally the results from transiesta and TBtrans calculations


